\chapter{Growth factors, transfer functions and linear input}\label{ch:growth-transfer}
\chapterstatus{written}{2026-11-14}

Initial conditions for simulations, the time stepping of COLA and LPT, and many analytic predictions need two things
from linear theory: the growth factors of Lagrangian perturbation theory as functions of time (and of scale when
gravity or the matter content makes growth scale dependent), and the transfer functions of each species at the
initial time, in the right gauge. This chapter describes the one growth solver of the library and the transfer
providers that read CAMB and CLASS output or run the Einstein--Boltzmann solver of
\cref{ch:background,ch:recombination,ch:linear-perturbations}.

\begin{modulebox}{Cosmology/Growth}
\file{Cosmology/Growth/LPTGrowthFactors.h} (namespace \nsp{FML::COSMOLOGY::GROWTH}; header-only): \cls{LPTGrowthEquations},
\cls{LPTGrowthSettings}, \cls{LPTGrowthFactors}, \code{radiation_matter_growing_mode}. Used by SimulationCode's
\code{GravityModel} and by \code{LPT::compute_LPT_growth_factors_LCDM}/\code{_GeffGLCDM}; \code{history()} gives a
\cls{GrowthHistory} (\file{LPT/GrowthHistory.h}) for the $N$-body time stepping.\par
Documentation: \file{Cosmology/Growth/README.md}. Tests: \file{Cosmology/Growth/tests/test_growth.cpp}.
\end{modulebox}
\begin{modulebox}{TransferProviders and InitialConditions}
\file{InitialConditions/TransferProvider.h} (the interface: species, gauges, units, coverage),
\file{InitialConditions/InternalTransferProvider.h}, \file{LinearInput.h}, \file{SpeciesIC.h};
\file{TransferProviders/MakeTransferProvider.h}, \file{TabulatedTransferProvider.h},
\file{CAMBFileTransferProvider.h}, \file{CLASSTransferProvider.h}, \file{GaugeView.h}, \file{TransferTablesFile.h};
scripts \file{TransferProviders/scripts/make_class_transfer_tables.py},
\file{SimulationCode/scripts/make_camb_transfer_tables.py}.\par
Documentation: \file{TransferProviders/README.md}, \file{InitialConditions/README.md},
\file{SimulationCode/INITIAL_CONDITIONS.md}. Tests: \file{TransferProviders/tests}, \file{tests/test_transfers.cpp},
\file{tests/test_camb.cpp}, \file{Cosmology/EinsteinBoltzmannSolver/tests/test_transfer_provider.cpp}.
\end{modulebox}

\section{Linear growth of structure}\label{sec:growth-linear}

Inside the horizon and after the radiation era, the clustering matter (CDM and baryons, or all matter) behaves as a
pressureless fluid with
\begin{equation}
\dot\delta+\frac1a\nabla\cdot[(1+\delta)\bm v]=0,\qquad
\dot{\bm v}+H\bm v+\frac1a(\bm v\cdot\nabla)\bm v=-\frac1a\nabla\Psi,\qquad
\nabla^2\Psi=4\pi Ga^2\mu\,\bar\rho_m\delta,
\end{equation}
with a possibly scale- and time-dependent effective gravitational constant $\mu=G_{\rm eff}/G$
(\cref{sec:ext-quasistatic}). Linearising and changing to $x=\ln a$ with $'=\dd/\dd x$ gives the linear growth
equation for $\delta(\vk,a)=D_1(k,a)\delta(\vk,a_{\rm ini})/D_1(k,a_{\rm ini})$,
\begin{equation}
D_1''+\left(2+\frac{\dd\ln E}{\dd x}\right)D_1'=\frac32\Om_m(x)\,\mu(k,x)\,D_1,
\label{eq:growth-D1}
\end{equation}
where $\Om_m(x)=\rho_m(x)/(\rho_{c0}E^2)$. In Einstein--de Sitter ($\Om_m=1$, $\dd\ln E/\dd x=-3/2$, $\mu=1$) the
growing mode is $D_1\propto a$; in $\Lambda$CDM the growth slows when $\Lambda$ dominates and
$f\equiv\dd\ln D_1/\dd x\simeq\Om_m(x)^{0.55}$.

\subsection{Higher orders in Lagrangian perturbation theory}

In LPT the particle positions are $\vx=\vq+\bm\Psi(\vq,t)$, expanded as
$\bm\Psi=\bm\Psi^{(1)}+\bm\Psi^{(2)}+\bm\Psi^{(3)}+\ldots$ with each order a product of a growth factor and a spatial
field \cite{Bouchet1995}: $\nabla_q\cdot\bm\Psi^{(1)}=-D_1\delta_{\rm ini}$, $\nabla_q\cdot\bm\Psi^{(2)}=
\frac{D_2}{2}\sum_{i\ne j}(\phi^{(1)}_{,ii}\phi^{(1)}_{,jj}-\phi^{(1)2}_{,ij})$, and the third-order terms $D_{3a}$ (the
determinant of the Hessian of $\phi^{(1)}$) and $D_{3b}$ (the mixed first--second-order term); the spatial
kernels are described in \cref{ch:lpt}. The growth factors obey (with $S\equiv\frac32\Om_m(x)\mu(k,x)$)
\begin{align}
D_1''+\left(2+\tfrac{\dd\ln E}{\dd x}\right)D_1'&=S\,s_1\,D_1,\label{eq:growth-orders}\\
D_2''+\left(2+\tfrac{\dd\ln E}{\dd x}\right)D_2'&=S\,s_2\,(D_2-N\,D_1^2),\nonumber\\
D_{3a}''+\left(2+\tfrac{\dd\ln E}{\dd x}\right)D_{3a}'&=S\,s_{3a}\,(D_{3a}-2D_1^3),\nonumber\\
D_{3b}''+\left(2+\tfrac{\dd\ln E}{\dd x}\right)D_{3b}'&=S\,s_{3b}\,(D_{3b}+D_1^3-D_2D_1).\nonumber
\end{align}
In EdS the growing solutions are
\begin{equation}
D_2=-\tfrac37D_1^2,\qquad D_{3a}=-\tfrac13D_1^3,\qquad D_{3b}=\tfrac{5}{21}D_1^3 .
\end{equation}
The factors in \cref{eq:growth-orders} generalise GR:
\begin{itemize}
\item $N(x,k)$ is the \emph{second-order source} of the theory of gravity (\code{second_order_source} of
      \cls{QuasiStaticGravity}): 1 in GR and in theories whose scalar field equation has no nonlinear term; smaller
      than 1 for DGP and the cubic Galileon, whose Vainshtein term already acts at second order.
\item $s_i(x,k)$ are per-order source factors (default 1). With massive neutrinos the CDM+baryon fluid feels the
      gravity of only the clustering part of the matter; with $\rho_m$ including the neutrinos one uses e.g.\
      $s_1=1-f_\nu$ above the free-streaming scale, or the source computed from the neutrino transfer functions
      (as SimulationCode does).
\item A velocity friction $\Gamma(x)$ (\code{friction}, default 0) adds $\Gamma D_i'$ to the damping term at every
      order. It describes a drag in the equation of motion $\dd\bm v/\dd x=-\Gamma\bm v+\text{force}$, for example
      dark scattering (\cref{sec:ext-fluids}). In EdS with constant $\Gamma$ the growing mode is $D_1\propto a^p$ with
      $p=\frac12\left(-b+\sqrt{b^2+6}\right)$, $b=\frac12+\Gamma$, and $D_2=c_2D_1^2$ with
      $c_2=-\frac32/(Q_2-\frac32)$, $Q_2=4p^2+2bp$; these are among the tests.
\end{itemize}

\subsection{Initial conditions in a matter + radiation universe}

At early times radiation cannot be neglected in $E$ but does not cluster. With $y=a/a_{\rm eq}$, a universe of matter
and smooth radiation has the linear equation
\begin{equation}
(1+y)\frac{\dd^2D}{\dd\ln y^2}+\frac y2\frac{\dd D}{\dd\ln y}-\frac32y\,D=0,
\end{equation}
whose growing mode is the M\'esz\'aros solution $D_1=1+\frac32y$ \cite{Meszaros1974}: no growth deep in the radiation
era, $D_1\propto a$ in matter domination. The higher orders obey the same operator with sources, and
\code{radiation_matter_growing_mode} solves them as power series $D=\sum_nc_ny^n$ ($c_0=0$, convergent for $y<1$),
which tend to the EdS values for $y\to\infty$. The growth solver starts at \code{a_min} with these series, with
$y=r/(1-r)$ and $r=\rho_m/E^2$ evaluated on the actual background (so the start is consistent with whatever species
the background has), and with the EdS values once $y\ge0.5$ (or with \code{eds_initial_conditions}). With
\code{start_in_radiation_era} (default) \code{a_min} is moved to $a_{\rm eq}/10$ when radiation matters at
\code{a_min} but it is not deep in the radiation era ($0.1<a_{\min}/a_{\rm eq}<10^4$). An \code{initial_rate_1(x,k)}
can replace $\dd\ln D_1/\dd x$ at the start (e.g.\ the CDM+baryon growth rate from transfer functions with massive
neutrinos).

\begin{implementation}[Implementation: \code{LPTGrowthFactors}]
\begin{itemize}
\item \cls{LPTGrowthEquations} is a set of functions: $E$, $\dd\ln E/\dd x$, the clustering matter density
      $\rho_m(x)$ (defaults from a \cls{Background}: $\Om_m(x)E^2$, all matter including massive neutrinos), $\mu$ and
      $N$ (from a \cls{QuasiStaticGravity} or user functions), the source factors $s_i$, the friction and the initial
      rate. They are called only while solving, in parallel over $k$, so they must be thread safe.
\item The four second-order ODEs are integrated on \code{n_a} $=1000$ points in $a\in[\code{a_min},\code{a_max}]=
      [10^{-4},10]$ and normalised so that $D_1=1$ at \code{a_normalize} ($D_2$ by $D_1^2$ and $D_3$ by $D_1^3$ at the same
      time); the results and the rates $f_i=\dd\ln D_i/\dd x$ are splined in $x$.
\item With \code{scale_dependent} the equations are solved on \code{n_k} $=200$ log-spaced wavenumbers
      $k\in[10^{-4},100]\hMpc$ (in units of $H_0/c$, \code{koverH0}) and stored as 2D splines; outside the grid $k$ is
      clamped, and \code{D1_k0} etc.\ give the large-scale solution. $E$, $\dd\ln E/\dd x$ and $\rho_m$ are then
      tabulated once on \code{n_tabulate} $=4000$ points (relative error $\sim10^{-9}$), since with $n_k$ solves the
      background evaluations would otherwise dominate.
\end{itemize}
\end{implementation}

\begin{cppcode}[Growth factors in GR and in $f(R)$]
#include <FML/Cosmology/Growth/LPTGrowthFactors.h>
#include <FML/Cosmology/Perturbations/QuasiStaticGravity_HuSawicky.h>
using namespace FML::COSMOLOGY;

auto bg = BACKGROUND::make_background("lcdm", params);
GROWTH::LPTGrowthSettings s;
s.a_normalize = 1.0 / 50.0;                         // D1 = 1 at z = 49
GROWTH::LPTGrowthFactors gr(*bg, s);                // GR
double D2 = gr.D2(std::log(0.5));                   // ~ -3/7 D1^2

s.scale_dependent = true;                           // mu(k, a): solve on a grid in k
auto fr = std::make_shared<PERTURBATIONS::QuasiStaticGravity_HuSawicky>(*bg, 1e-5, 1.0);
GROWTH::LPTGrowthFactors grfr(*bg, s, fr);
double koverH0 = 1.0 * FML::UTILS::SI::c_over_H0_mpch;   // k = 1 h/Mpc
double boost = grfr.D1(0.0, koverH0) / gr.D1(0.0);
\end{cppcode}

\begin{validation}[Validation: growth factors]
\begin{itemize}
\item EdS: $D_1=a/a_{\rm norm}$, $D_2/D_1^2=-3/7$, $D_{3a}/D_1^3=-1/3$, $D_{3b}/D_1^3=5/21$ to $10^{-6}$, rates
      $f_1=1$, $f_2=2$, $f_{3a}=3$ to $10^{-5}$; with constant friction $\Gamma=0.7,-0.3$ the analytic $a^p$ solutions
      above to $10^{-6}$.
\item Matter + radiation: \code{a_min} moved to $a_{\rm eq}/10$, the M\'esz\'aros mode $1+\frac32y$ to $10^{-6}$, and
      $D_2\to-\frac37D_1^2$ for $a\gg a_{\rm eq}$.
\item A \cls{QuasiStaticGravity} with $\mu=1$ reproduces GR to $10^{-14}$; $f(R)$ at $k=10^{-4}\hMpc$ is GR to
      $\order{|f_{R0}|}$ and grows 30--70\% more at $100\hMpc$; for DGP $D_1$ does not depend on $N$ while $|D_2|$ is
      reduced by the second-order source.
\end{itemize}
\end{validation}

\section{Transfer functions}\label{sec:transfer}

A \cls{TransferProvider} gives the linear transfer functions of a set of species,
\code{Species} $\in$ \{\code{CDM}, \code{Baryon}, \code{CDMBaryon}, \code{MassiveNeutrino}, \code{TotalMatter}\}, as
functions of $k$ (in $\hMpc$) and the scale factor $a$:
\begin{equation}
T_s(k,a)\equiv\frac{\delta_s(k,a)}{\mathcal R(k)},\qquad
T^\theta_s(k,a)\equiv\frac{\theta_s(k,a)}{aH\,\mathcal R(k)},\qquad
f_s(k,a)=\frac{\dd\ln\delta_s}{\dd\ln a},
\end{equation}
per unit primordial adiabatic curvature in the CAMB sign convention, with $\theta=\nabla\cdot\bm u$, $\bm u=a\,\dd\vx/\dd t$
(\code{density_transfer}, \code{velocity_divergence_transfer}, \code{density_growth_rate}). The matter power
spectrum follows as
\begin{equation}
P_s(k,a)=\frac{2\pi^2}{k^3}\mathcal P_\mathcal R(k)\,T_s^2(k,a),\qquad
\mathcal P_\mathcal R=A_s\left(\frac{k}{k_{\rm pivot}}\right)^{n_s-1},
\end{equation}
in $(\Mpch)^3$. \code{density_over_k_mpc2} returns the CAMB representation $T/k^2$ in Mpc$^2$ (with the physical $k$ in
Mpc$^{-1}$). Composite species are weighted by the background densities $\rho_i(a)$ of the model (densities) and
$\rho_i+p_i$ (velocities). A density growth rate is never silently substituted for a velocity transfer: the direct
velocities are Newtonian-gauge quantities, while the growth rate is a derivative of the density in the provider's
gauge.

\subsection{Gauges}\label{sec:transfer-gauges}

Particle initial conditions with the Zel'dovich displacement need the density in the \emph{CDM-synchronous} gauge,
the frame comoving with the dark matter, where $\delta$ is what the particles' Lagrangian displacements produce.
\code{Gauge} $\in$ \{\code{Newtonian}, \code{SynchronousCDM}, \code{NBody}\}.

\paragraph*{Newtonian $\to$ CDM-synchronous} is the time shift that brings the CDM to rest
\cite{Ma1995}, applied with the same shift to every species $s$ (equation of state $w_s=p_s/\rho_s$):
\begin{equation}
\delta_s^{\rm sync}=\delta_s^N+3(1+w_s)\left(\frac{\Hp}{ck}\right)^2\frac{\theta_{\rm cdm}^N}{\Hp},\qquad
\frac{\Hp}{ck}=\frac{aE(a)}{2997.92458\,k[\hMpc]},
\end{equation}
which requires the Newtonian CDM velocity. The velocities are passed through unchanged, so a converted provider has
CDM-synchronous densities and Newtonian velocities. The opposite direction, synchronous $\to$ Newtonian, needs the
metric combination $(h'+6\eta')/2k^2$ and is not supported. On the largest scales the conversion cancels terms $10^5$
times the result, so the converted density is accurate to $3\times10^{-3}$ at $k=10^{-4}\hMpc$ and to $7\times10^{-4}$
above $10^{-3}\hMpc$.

\paragraph*{The $N$-body gauge.} In this gauge \cite{Fidler2015} the CDM+baryon particles follow the Newtonian
equations of motion to linear order in GR, with an extra source $\delta_{\rm extra}$ in the Poisson equation,
\begin{equation}
\delta_{cb}''+\Hp\delta_{cb}'=4\pi Ga^2\rho_{cb}\,(\delta_{cb}+\delta_{\rm extra})\qquad(\text{conformal time}),
\end{equation}
which carries the effect of radiation and of the relativistic terms on large scales. The internal provider computes
the $N$-body gauge densities $\delta_s=\delta_s^N+(1+w_s)S$, with the shift to the comoving density
$S=3\Hp\theta_{\rm tot}/(ck)^2=3(\Psi-\Phi')/D$, $D=1-\dd\ln\Hp/\dd x-\Om_K(x)$, and the source
$\delta_{\rm extra}=-(\Theta'+\Hp\Theta)/(4\pi Ga^2\rho_{cb})-\delta_{cb}$, where $\Theta=\theta^N_{cb}+H_T'$ is the
$N$-body gauge velocity and $H_T'=(3\Phi-S)'$. To avoid the cancellation of two large terms on small scales, $(ck)^2\Psi$
is rewritten with the comoving Poisson equation so that $\delta_{\rm extra}$ becomes a sum of small terms (photons,
neutrinos and dark energy), and the $x$-derivatives of $3\Phi-S$ and of the baryon--CDM velocity difference are local
least-squares fits on the solver's grid. SimulationCode uses these tables with \code{force_nbody_gauge}.

\subsection{The providers}

\code{make_transfer_provider(settings)} (\file{TransferProviders/MakeTransferProvider.h}) builds a provider from one of
three sources; the user declares the gauge of the input and asks for the gauge of the output:

\begin{center}\small
\begin{tabularx}{\textwidth}{@{}lXX@{}}
\toprule
\code{source} & Input & Gauge handling \\
\midrule
\code{camb} & A CAMB transfer info file (\cls{CAMBFileTransferProvider}): \code{<folder> <n>}, then
  \code{<file> <z>} per line; columns by name from the header (or by position: 13 CAMB, 9 AxionCAMB, 7 old CAMB) &
  densities are CDM-synchronous; CDM and baryon velocities Newtonian; no neutrino velocity \\
\code{class} & Native CLASS output (\code{format = class}, \code{output = mTk[,vTk]}, \code{write_background = yes});
  \cls{CLASSTransferProvider} reads columns by name and the redshift from the header & declared
  \code{synchronous} or \code{newtonian}; Newtonian is converted at the tabulated redshifts (needs \code{vTk}) \\
\code{internal} & FML's Einstein--Boltzmann solver through \cls{InternalTransferProvider}: any model of the registry,
  curvature, massive neutrinos & both gauges (and the $N$-body gauge) from one solve \\
\bottomrule
\end{tabularx}
\end{center}

\paragraph*{Sign and unit conventions.} All providers return $\delta/\mathcal R$ and $\theta/(aH\mathcal R)$ in the CAMB
sign convention. CLASS normalises to unit initial curvature with the opposite sign, so the CLASS provider returns
$-d_i$ and $-t_i/(aH)$; CAMB's columns are $T=\delta/k^2$ in Mpc$^2$ and $-\theta/(aHk^2)$, so the CAMB provider
returns $Tk^2$ and $-T_vk^2$. \code{TotalMatter} is the $\rho$-weighted sum of CDM, baryons and massive neutrinos in the
provider's gauge (CLASS's \code{d_m}, the comoving matter density, is not used). The declared gauge is checked as far
as the columns allow: a synchronous CLASS run never has \code{t_cdm}, a Newtonian run with \code{vTk} has it.

\paragraph*{Interpolation and extrapolation.} The tabulated providers share \cls{TabulatedTransferProvider}:
bicubic interpolation in $(\ln a,\ln k)$ with four or more redshifts (bilinear with two or three), growth rates from
the derivative of a 4-point Lagrange polynomial in $\ln a$ through the tables (so CAMB and CLASS give the same growth
rates from the same redshift sampling), and a \code{RangePolicy}: \code{Reject} (outside the tables is an error) or
\code{SourceExtrapolation} (default; a power law in $k$ above $k_{\max}$ through the last $0.1$ in $\ln k$, and the end
values elsewhere). The tables must be finely sampled in redshift and extend past the redshifts where growth rates are
used: with 7 redshifts between $z=49$ and 0 the growth rate was off by up to 29\% at the top of the table; with 40
redshifts up to $z=60$ CAMB and CLASS agree to $2\times10^{-4}$. The scripts' default (60 redshifts uniform in
$\ln(1+z)$ up to $z=200$) is safe.

\paragraph*{The internal provider.} It runs the solver when it is made (a few minutes with $k\in[10^{-5},10]\hMpc$ and
400 modes) and owns the solution. The perturbations are solved to $a=1.1$, because a table ending at $a=1$ distorted
the growth rate at $z=0$ by $2.6\times10^{-3}$ through the spline end; growth rates are centred differences
($\pm10^{-3}$ in $\ln a$). Massive neutrinos are \code{deg_ncdm} degenerate species whose mass gives
$\Om_{\rm ncdm}=$ \code{OmegaMNu}, or any masses with \code{m_nu_list}; the other gauge views share the solution
(\code{with_density_gauge}). With \code{axion_as_cdm} (\code{model = axion}) the CDM species is CDM plus the
ultralight axion, for simulations whose particles are all the dark matter. Isocurvature modes and running are
rejected. Near the solver's massive-neutrino fluid switch ($k\tau=31$) the neutrino velocities differ from CLASS at
high redshift (40\% at $k=0.03\hMpc$, $z=49$; $10^{-3}$ at $z\le2$); for neutrino particles started early, raise
\code{ncdm_fluid_trigger}.

\paragraph*{Transfer tables.} \code{tabulate_transfers(provider, a, k, A_s, n_s, k_pivot, policy)} tabulates any
provider, and \cls{TransferTablesProvider} makes a provider of the tables (cheap to evaluate, and extrapolated like
CAMB and CLASS tables). \file{TransferTablesFile.h} writes and reads them as one self-describing text file: \code{\#
key = value} header lines (cosmology, primordial spectrum, gauges, solver settings) and one row per $(a,k)$ with
\code{d_<s>} $=\delta/\mathcal R$ and \code{t_<s>} $=\theta/(aH\mathcal R)$ for $s=$ \code{cdm}, \code{b}, \code{cb},
\code{ncdm}, \code{m}, and the $N$-body gauge source. SimulationCode writes them with
\code{./nbody --write-transfers params.lua tables.txt} and reads them with \code{ic_type_of_input = "tables"}, checking
the cosmology and gauge against the simulation, so the solver runs once per cosmology.

\paragraph*{Gauge views.} \cls{GaugeView} wraps any provider that has the Newtonian CDM velocity and applies the
Newtonian$\to$CDM-synchronous conversion, for sources that cannot convert at their own tabulated redshifts.

\begin{cppcode}[A transfer provider from CLASS output in the Newtonian gauge]
#include <FML/TransferProviders/MakeTransferProvider.h>
using namespace FML::INITIALCONDITIONS;

TransferSourceSettings s;
s.source       = "class";
s.class_root   = "output/run_00_";        // run_00_z1_tk.dat, ..., run_00_background.dat
s.input_gauge  = Gauge::Newtonian;        // what CLASS was run with
s.output_gauge = Gauge::SynchronousCDM;   // what particle ICs need (the default)
s.A_s = 2.1e-9; s.n_s = 0.965;            // k_pivot_mpc = 0.05 by default
auto provider = make_transfer_provider(s);
double T = provider->density_transfer(Species::CDMBaryon, 0.1, 1.0 / 50.0);     // k [h/Mpc], a
double f = provider->density_growth_rate(Species::CDMBaryon, 0.1, 1.0 / 50.0);
\end{cppcode}

\begin{validation}[Validation: transfer providers (CLASS 3.4.1, CAMB 2.0.5, 0.1\,eV neutrino, 40 redshifts to $z=60$)]
\begin{center}\small
\begin{tabularx}{\linewidth}{@{}Xr@{}}
\toprule
Check & Agreement \\
\midrule
Newtonian CLASS run converted vs synchronous run, CDM/b/cb/total, $10^{-3}\le k\le1\hMpc$ & $7\times10^{-4}$ \\
The same at $k=10^{-4}\hMpc$ (unconverted: $5\times10^2$) & $3\times10^{-3}$ \\
CLASS vs CAMB, CDM/b/cb/total; $P_{cb}(z=0)$; $f_{cb}$ & $5\times10^{-4}$; $8\times10^{-4}$; $2\times10^{-4}$ \\
Internal vs CLASS, both gauges, massless $\nu$ ($0.1\,$eV $\nu$), $10^{-3}\le k\le1$ & $7\times10^{-4}$ ($9\times10^{-4}$) \\
Internal vs CLASS: $f_{cb}$; Newtonian CDM velocity & $8\times10^{-4}$; $6\times10^{-4}$ \\
Internal: \code{GaugeView} on its Newtonian view vs its own synchronous view & $10^{-5}$ \\
CAMB reader vs the former \code{LinearTransferData} on the same files & $2\times10^{-5}$ \\
\bottomrule
\end{tabularx}
\end{center}
Massive neutrinos below the free-streaming scale depend on the codes' approximations: with the scripts' defaults
(CLASS without its neutrino fluid approximation, CAMB \code{AccuracyBoost = 3}) $\delta_\nu$ agrees to 1.4\% at $z=0$ for
$k\le1\hMpc$, but differs by tens of per cent at $z=49$ where it is tiny.
\end{validation}

\section{Preparing the linear input of initial conditions}\label{sec:linear-input}

\paragraph*{Simple inputs.} \file{LinearInput.h} prepares the CDM+baryon $P(k)$, the primordial spectrum and
$T(k)$ at the initial redshift from a two-column $P(k)$ file ($k$ in $\hMpc$, $P$ in $(\Mpch)^3$), a single transfer
function ($T$ in Mpc$^2$) or a provider, without a Boltzmann solve. The arrays are generated at the input redshift
and brought to the initial time with the growth factors of the simulation's gravity model, so choosing
\code{ic_input_redshift} selects between \emph{back-scaling} (input at $z=0$, scaled back with $D_1$: the initial
conditions reproduce the linear $P(k)$ of the target model today in the simulation's own dynamics) and forward
initial conditions (input at $z_{\rm ini}$: the actual linear state at that time, including e.g.\ radiation effects).

\paragraph*{Two species.} For separate CDM and baryon particles (\file{SpeciesIC.h}) the two fields share one random
realisation, $\delta_i(\vk)=\delta_{cb}(\vk)R_i(k,z_{\rm ini})$ with $R_i=T_i/T_{cb}$ ($w_cR_c+w_bR_b=1$). The 1LPT
velocity uses $f_i=f_{cb}+\dd\ln R_i/\dd\ln a$, where $f_{cb}$ is the growth rate of the gravity model, so that the
total CDM+baryon velocity is exactly the growing mode and the relative velocity comes from the transfer functions.
With \code{particle_baryon_relative_mode = mass} the relative mode is instead put in the particle masses,
$\epsilon_i(k)=[T_i(k,0)-T_{cb}(k,0)]/T_{cb}(k,z_{\rm ini})$ ($w_c\epsilon_c+w_b\epsilon_b=0$), following Hahn, Rampf and
Uhlemann (\arxiv{2008.09124}): both species then follow the total CDM+baryon flow, and the relative mode, constant
in linear theory, is exact today. \code{mass_ini} uses $\epsilon_i=T_i/T_{cb}-1$ at $z_{\rm ini}$ with separate
velocities, so the relative mode grows from $z_{\rm ini}$; the default \code{displacement} displaces each species
by its own linear field.

\paragraph*{Glass loads.} Initial conditions on a glass instead of a lattice (GlassIC, \cref{ch:lpt}) sample the
same provider-based Fourier fields at the irregular particle positions.

\begin{paramfile}[SimulationCode: initial conditions from the internal solver or from tables]
ic_type_of_input  = "internal"            -- or "camb", "class", "tables"
ic_input_filename = ""                    -- optional solver .ini (precision, recombination)
ic_internal_solver_parameters = {"Yp=0.245", "precision=fast"}
ic_input_redshift = ic_initial_redshift   -- transfers at zini, no back-scaling
-- once per cosmology:  ./nbody --write-transfers parameterfile.lua transfers_lcdm.txt
-- then:                ic_type_of_input = "tables"; ic_input_filename = "transfers_lcdm.txt"
\end{paramfile}
With \code{internal}, the solver runs on task 0 with the simulation's cosmology and model (background, massive
neutrinos as three degenerate species of mass $M_\nu/3$, \code{cosmology_model}, primordial spectrum) and sends its
tables to the other tasks; $k$ is sampled with 15 points per decade from $10^{-4}\hMpc$ and 50 per decade in
$[0.01,0.6]\hMpc$ (the BAO) up to $\max(5\hMpc,1.5\times$ the largest mode of the IC grid), and the tables cover 128
scale factors from $a_{\rm ini}/2$ to $1.1$ ($\sim30$\,s on two cores). Simulations started from these tables agree
with CAMB-input runs to $3\times10^{-4}$ in $P(k)$ at $z=0,1$ for $k<0.5\hMpc$ without neutrinos and to $0.4\%$ with
0.4\,eV. The details of the initial-condition generation are in \cref{ch:lpt,ch:simulation-code}.
